% Normal records and common acceptance rules; full encodings are in Appendix A.
\subsection{회계 기록과 상태}
\label{sec:objects}

원장은 상세 회계 값을 숨기되 각 객체의 생성·소비 여부를 관리한다.
Note는 다음 생산·이전에 쓰고, Voucher는 인수 또는 회수로 해결한다.
발급된 Claim은 공개 결과이므로 다시 소비하지 않는다.

회계 값은 질량 $q$, 재활용 크레딧 $a$, 탄소 $e$이다.
모두 음이 아닌 정수이고 $0\leq a\leq q$를 만족한다.
$q,a$는 같은 질량 단위·배율, $e$는 고정 배율의
$\mathrm{kgCO_2e}$를 쓴다. $a$는 배분한 회계 근거이지 실제 재활용
함량의 측정값은 아니다. $e$는 수용한 초기 값과 이후 증분의 배분값이다.

참여자는 원래 값과 비밀 난수인 개방값을 보관하고, 원장에는
이를 숨기면서 나중에 다른 값으로 바꾸기 어렵게 하는
커미트먼트 $cm$을 등록한다. 해시 $H$는 입력에서 고정 길이의 값을 만든다.
회계 증표에는 회계 값 외에 문서 해시 $d$, 소유자 $o$와
역할 $\rho$가 담긴다. 역할은 일반용($\eligible$) 또는
폐기물($\waste$)이다. 폐기물 증표는 $a=e=0$이고
종료(Exit)로만 소비한다.

같은 기록의 재사용을 막기 위해 Note의 소비 식별자를 계산한다.
\begin{equation}
 nf=H_{\mathsf{Nullifier}}(cm,sk)
 \label{eq:nf}
\end{equation}
생성 때 암호화해 남긴 식별자를 소비 때 공개한다.
원장은 재사용을 거부하고, 감사자는 두 기록을 연결할 수 있다.

Voucher의 공개 커미트먼트는 $rv$이며 송·수신자와 회수 기한은
개방값에 담긴다. 인수(Proceed)와 회수(Recall)는 다른 권한으로
같은 해결 식별자 $rvnf$를 소비하므로 한 경로만 수락된다.
정확한 계산은 부록~\ref{app:formats}에 둔다.

제품 주장에는 제품명·로트 식별자·단위를 해시한 $d$,
발급 정책 $p$, 발급을 구분하는 공개 보조값인 논스 $c$를 연결한다.
\begin{equation}
 h=H_{\mathsf{Issue}}(d,p,c),\qquad \mathsf{DPPClaim}=(p,h,c)
 \label{eq:claim}
\end{equation}
Claim은 종단이므로 소유자·개방값·미래 소비 식별자·사용 상태가 없다.
객체 참조는
$\mathsf{NoteRef}(cm)$, $\mathsf{VoucherRef}(rv)$,
$\mathsf{ClaimRef}(h)$처럼 타입을 구분한다.
숫자가 같아도 서로 다른 타입의 객체를 합치지 않는다.

원장은 등록·소비·사용 상태를 구분해 관리한다.
Note와 Voucher의 등록은 머클 트리의 포함 증명으로 확인한다.
머클 트리는 기록 목록을 하나의 해시값으로 요약해 특정 기록의
포함 여부를 증명하는 구조다.
타입별 $\spend[\nu]$는 미소비이면 0, 소비됐으면 그 이벤트
식별자 $\aid$이며 한 번 기록한 소비는 바꾸지 않는다.
각 출력의 생성 이벤트도 $\mathsf{producerOf}$로 고정한다.

사용 상태는 활성($\stActive$), 동결($\stFrozen$),
사용 금지($\stRevoked$) 중 하나다.
\begin{equation}
 \begin{gathered}
 \stActive\longrightarrow\stFrozen,\\
 \stFrozen\longrightarrow\{\stActive,\stRevoked\}
 \end{gathered}
 \label{eq:status}
\end{equation}
상태 관리자만 미소비 객체의 상태를 바꾼다. 종료를 포함한
\textbf{모든 소비는 실행 시점에 미소비·활성인 입력만 허용한다}.
과거의 등록 증명만으로 현재 사용을 허용하지 않는다.

\subsection{거래 검증과 감사 자료}
\label{sec:transitions}
\label{sec:binding-interface}

거래를 수락하려면 회계 계산, 실제 입력·출력과 감사 자료의 대응,
현재 입력의 사용 가능 상태를 모두 검사한다. 원료 하나를 둘로
나눴다면 감사 자료에도 바로 그 입력과 두 출력이 반영되어야 한다.

감사 기록(AuditRecord)은 거래 종류·정책·출력 목록과 함께
실제 입력인 \emph{부모}의 참조와 출력의 미래 소비 식별자를
암호화한다. 부모 참조 목록 뒤에 미래 식별자 목록을 이어 붙여
하나의 벡터, 즉 순서 있는 값 목록으로 암호화한다.
영지식증명은 회계 계산과 이 목록의 암호화가 같은 실제 입력·출력에서
계산됐는지 함께 검사한다.
개방값을 $z$, 거래의 비공개 입력 전체를 $w$라 하면,
핵심 연결은 다음 세 식이다.
\begin{equation}
 \begin{aligned}
 r_i^{\rm out}&=\mathsf{Ref}(z_i^{\rm out}),\\
 \nu_i^{\rm in}&=\mathsf{SpendId}(z_i^{\rm in},sk_i),\\
 ct&=\mathsf{AuditEnc}(pk_{\mathsf{audit}},M(w);\eta)
 \end{aligned}
 \label{eq:audit-binding-interface}
\end{equation}
세 식은 각각 출력 참조, 입력의 소비 식별자, 실제 자료에서
계산한 감사 암호문을 뜻한다.
$M(w)$는 부모의 원시 참조와 출력 미래 식별자를 정해진 순서로
모은 목록이고, $\eta$는 이 거래의 암호화 난수다.
예를 들어 이전 거래에서는 입력 Note의 $cm$ 뒤에
출력 Voucher의 $rvnf$와 잔여 Note의 $nf$를 놓는다.
각 자리의 타입·용도와 목록 길이는 거래 회로가 정하고,
암호문 전체를 회계 계산·입출력 참조와 같은 증명에 결합한다.
키는 승인된 $PK_B$로 고정하며 증명자가 임의로 바꿀 수 없다.
암호 계산에는 거래별 문맥을 따로 넣지 않는다.
서로 다른 거래에는 새로운 독립 난수를 사용한다.

적법한 소비에는 생성 때 암호화한 식별자를 써야 한다.
Note는 개방값의 결합성과 소유권 비밀의 유일한 표현을,
Voucher는 두 해결 경로의 공통 식별자를 요구한다.
이를 어기면 같은 기록을 다른 식별자로 다시 쓸 수 있다.
세부 조건과 반례는 부록~\ref{app:formats}에 둔다.
